\documentclass{article}
\usepackage{amsmath}
\usepackage{amssymb}
\usepackage{geometry}
\usepackage{graphicx}
\usepackage{xepersian}
\settextfont{XB Niloofar}

\geometry{a4paper, margin=1in}

\title{\lr{Maximum Likelihood Estimation (MLE)} \\ (تخمین بیشینه درست‌نمایی)}
\author{}
\date{}
\begin{figure}[h!]
  \centering
  \includegraphics[width=0.75\textwidth]{img_page_2_1.png}
  \caption{تخمین درست‌نمایی بیشینه (MLE) برای پارامترهای توزیع}
\end{figure}




\begin{document}

\maketitle

\section*{فهرست مطالب}
\begin{enumerate}
    \item مقدمه
    \item تابع درست‌نمایی (\lr{The Likelihood Function})
    \item استخراج ریاضی: گاوسی تک‌متغیره (\lr{Univariate Gaussian})
    \item تعمیم به گاوسی چندمتغیره (\lr{Multivariate Gaussian})
    \item ویژگی‌ها و محدودیت‌ها
    \item ارتباط با دسته‌بندی بیز (\lr{Bayesian Classification})
    \item پل ارتباطی به موضوعات پیشرفته
    \item خلاصه
\end{enumerate}

\newpage

\section*{۱. مقدمه}
در دسته‌بندی الگوهای آماری، یک فرض رایج این است که داده‌ها توسط یک توزیع احتمال (\lr{probability distribution}) زمینه‌ای با پارامترهای ناشناخته تولید می‌شوند. به عنوان مثال، ممکن است فرض کنیم ویژگی‌های $x$ از یک توزیع گاوسی (نرمال) پیروی می‌کنند. برای اتخاذ تصمیمات بهینه (مانند دسته‌بندی با استفاده از تئوری تصمیم‌گیری بیز)، باید پارامترهای این توزیع‌ها را از یک مجموعه‌داده آموزشی تخمین بزنیم.

تخمین بیشینه درست‌نمایی (\lr{Maximum Likelihood Estimation} یا \lr{MLE}) یک روش پارامتریک استاندارد است که برای تخمین پارامترهای ناشناخته $\theta$ از یک تابع چگالی احتمال $p(x|\theta)$ با داشتن مجموعه‌داده $D$ استفاده می‌شود. اصل کار انتخاب مقادیر پارامتری است که داده‌های مشاهده‌شده را محتمل‌تر کند.

\section*{۲. تابع درست‌نمایی}
تابع درست‌نمایی که با $L(\theta|D)$ نشان داده می‌شود، احتمال مشاهده مجموعه‌داده خاص $D$ را به عنوان تابعی از پارامترهای $\theta$ نشان می‌دهد.

\subsection*{تمایز کلیدی}
تشخیص تفاوت بین تابع چگالی احتمال $p(x|\theta)$ و تابع درست‌نمایی $L(\theta|D)$ بسیار حیاتی است:
\begin{enumerate}
    \item $p(x|\theta)$: توزیع داده‌های $x$ را برای پارامترهای ثابت $\theta$ توصیف می‌کند.
    \item $L(\theta|D)$: احتمال مجموعه‌داده $D$ را به ازای داده‌های مشاهده‌شده ثابت، به عنوان تابعی از $\theta$ توصیف می‌کند.
\end{enumerate}
تابع درست‌نمایی در فضایی متفاوت از تابع چگالی احتمال قرار دارد. فرم عملکردی درست‌نمایی لزوماً با فرم عملکردی چگالی یکسان نیست.

مقداری که درست‌نمایی را بیشینه می‌کند با $\hat{\theta}$ نشان داده می‌شود؛ این مقدار لگاریتم درست‌نمایی (\lr{log-likelihood}) یعنی $l(\theta)$ را نیز بیشینه می‌کند.

\section*{۳. استخراج ریاضی: گاوسی تک‌متغیره}
با ساده‌ترین حالت شروع می‌کنیم: تخمین پارامترهای یک توزیع گاوسی تک‌متغیره. این کار به عنوان پایه‌ای برای درک حالت چندمتغیره مورد استفاده در دسته‌بندی الگو عمل می‌کند.

\subsection*{تعریف مسئله}
فرض کنید $D = \{x_1, x_2, \dots, x_n\}$ مجموعه‌داده‌ای شامل $n$ نمونه مستقل و هم‌توزیع (\lr{independent and identically distributed} یا \lr{i.i.d.}) باشد. فرض می‌کنیم هر نمونه از یک توزیع گاوسی تک‌متغیره استخراج شده است:
\[
p(x|\mu, \sigma^2) = \frac{1}{\sqrt{2\pi\sigma^2}} \exp\left( -\frac{(x-\mu)^2}{2\sigma^2} \right)
\]
که در آن $\mu$ میانگین و $\sigma^2$ واریانس است.

\subsection*{تابع درست‌نمایی}
تابع درست‌نمایی $L(\mu, \sigma^2)$ با فرض استقلال به صورت زیر است:
\[
L(\mu, \sigma^2) = \prod_{k=1}^n p(x_k|\mu, \sigma^2) = \left( \frac{1}{\sqrt{2\pi\sigma^2}} \right)^n \exp\left( -\sum_{k=1}^n \frac{(x_k-\mu)^2}{2\sigma^2} \right)
\]

\subsection*{تابع لگاریتم درست‌نمایی}
بیشینه‌سازی مستقیم $L$ به دلیل حاصل‌ضرب توابع نمایی دشوار است. از آنجا که لگاریتم یک تابع اکیداً صعودی است، بیشینه‌سازی لگاریتم درست‌نمایی $l(\mu, \sigma^2) = \log L(\mu, \sigma^2)$ همان تخمین‌ها را به دست می‌دهد:
\[
l(\mu, \sigma^2) = -\frac{n}{2} \log(2\pi) - \frac{n}{2} \log(\sigma^2) - \frac{1}{2\sigma^2} \sum_{k=1}^n (x_k-\mu)^2
\]

\subsection*{تخمین میانگین ($\mu$)}
با گرفتن مشتق جزئی نسبت به $\mu$ و قرار دادن آن برابر صفر:
\[
\frac{\partial l}{\partial \mu} = \frac{1}{\sigma^2} \sum_{k=1}^n (x_k - \mu) = 0 \implies \sum_{k=1}^n x_k - n\hat{\mu} = 0 \implies \hat{\mu} = \frac{1}{n} \sum_{k=1}^n x_k
\]
نتیجه: تخمین‌گر \lr{MLE} برای میانگین همان میانگین نمونه است.

\subsection*{تخمین واریانس ($\sigma^2$)}
با گرفتن مشتق جزئی نسبت به $\sigma^2$ (با فرض $\tau = \sigma^2$) و قرار دادن آن برابر صفر:
\[
\frac{\partial l}{\partial \tau} = -\frac{n}{2\tau} + \frac{1}{2\tau^2} \sum_{k=1}^n (x_k - \mu)^2 = 0 \implies -n\tau + \sum_{k=1}^n (x_k - \mu)^2 = 0
\]
با جایگذاری میانگین نمونه $\hat{\mu}$ به جای $\mu$:
\[
\hat{\sigma}^2 = \frac{1}{n} \sum_{k=1}^n (x_k - \hat{\mu})^2
\]
نتیجه: تخمین‌گر \lr{MLE} برای واریانس همان واریانس نمونه (با مخرج $n$) است.

\section*{۴. تعمیم به گاوسی چندمتغیره}
در دسته‌بندی الگو، ویژگی‌ها معمولاً بردار هستند ($x \in \mathbb{R}^d$). تابع چگالی گاوسی چندمتغیره به صورت زیر است:
\[
p(x|\mu, \Sigma) = \frac{1}{(2\pi)^{d/2} |\Sigma|^{1/2}} \exp\left( -\frac{1}{2} (x-\mu)^T \Sigma^{-1} (x-\mu) \right)
\]
تخمین‌های \lr{MLE} با به کارگیری همان روش بیشینه‌سازی لگاریتم درست‌نمایی به دست می‌آید:
\begin{itemize}
    \item \textbf{بردار میانگین ($\mu$):} 
    \[
    \hat{\mu} = \frac{1}{n} \sum_{k=1}^n x_k
    \]
    \item \textbf{ماتریس کوواریانس ($\Sigma$):} 
    \[
    \hat{\Sigma} = \frac{1}{n} \sum_{k=1}^n (x_k - \hat{\mu}) (x_k - \hat{\mu})^T
    \]
\end{itemize}

\section*{۵. ویژگی‌ها و محدودیت‌ها}
\begin{enumerate}
    \item \textbf{انحراف (\lr{Bias}):}
    \begin{itemize}
        \item تخمین برای میانگین $\hat{\mu}$ بدون انحراف (\lr{unbiased}) است.
        \item تخمین برای واریانس $\hat{\sigma}^2$ (و $\hat{\Sigma}$) دارای انحراف (\lr{biased}) است. این روش به طور سیستماتیک واریانس واقعی را کمتر از مقدار واقعی تخمین می‌زند زیرا از مخرج $n$ استفاده می‌کند. تخمین‌گر بدون انحراف برابر با $s^2 = \frac{1}{n-1} \sum (x_k - \hat{\mu})^2$ است.
    \end{itemize}
    \item \textbf{سازگاری (\lr{Consistency}):} با میل کردن تعداد نمونه‌ها به سمت بی‌نهایت ($n \to \infty$)، تخمین‌های \lr{MLE} به مقادیر واقعی پارامترها همگرا می‌شوند.
    \item \textbf{ناوردا بودن (\lr{Invariance}):} اگر $\hat{\theta}$ تخمین‌گر \lr{MLE} برای $\theta$ باشد، برای هر تابع $g(\theta)$، تخمین‌گر \lr{MLE} آن برابر با $g(\hat{\theta})$ است.
    \item \textbf{حساسیت به داده‌های پرت (\lr{Sensitivity to Outliers}):} داده‌های پرت (\lr{outliers}) می‌توانند به طور قابل توجهی تخمین‌های میانگین و کوواریانس را منحرف کنند.
\end{enumerate}

\section*{۶. ارتباط با دسته‌بندی بیز}
نقش اصلی \lr{MLE} امکان‌پذیر ساختن دسته‌بندی بیز (\lr{Bayesian Classification}) در صورت ناشناخته بودن پارامترها است.

\subsection*{رویکرد جایگذاری (\lr{Plug-In Approach})}
\begin{enumerate}
    \item \textbf{آموزش:} برای هر کلاس $\omega_i$، از داده‌های آموزشی $D_i$ برای محاسبه تخمین‌های \lr{MLE} یعنی $\hat{\mu}_i$ و $\hat{\Sigma}_i$ استفاده کنید.
    \item \textbf{ساخت:} چگالی شرطی کلاس تخمین‌زده شده را تشکیل دهید: $\hat{p}(x|\omega_i) = N(x|\hat{\mu}_i, \hat{\Sigma}_i)$.
    \item \textbf{تصمیم‌گیری:} قاعده تصمیم بیز را با استفاده از این چگالی‌های تخمینی اعمال کنید. این روش به عنوان دسته‌بندی‌کننده جایگذاری (\lr{Plug-In Classifier}) شناخته می‌شود.
\end{enumerate}

\subsection*{تاثیر بر مرزهای تصمیم‌گیری}
\begin{itemize}
    \item \textbf{تحلیل ممیز خطی (\lr{LDA}):} فرض می‌کند همه کلاس‌ها دارای ماتریس کوواریانس مشترک $\Sigma$ هستند. مرز تصمیم خطی (ابرصفحه) می‌شود.
    \item \textbf{تحلیل ممیز درجه دوم (\lr{QDA}):} فرض می‌کند هر کلاس کوواریانس منحصربه‌فرد $\Sigma_i$ دارد. مرز تصمیم درجه دوم (بیضی‌گون با ابعاد بالا) می‌شود.
\end{itemize}

\section*{۷. پل ارتباطی به موضوعات پیشرفته}
\begin{itemize}
    \item \textbf{تحلیل مولفه‌های اصلی (\lr{PCA}):} می‌تواند به عنوان یک مسئله بیشینه درست‌نمایی برای مدل گاوسی با محدودیت‌ها مشتق شود.
    \item \textbf{مدل‌های مخلوط گاوسی (\lr{GMM}):} گاوسی‌های منفرد را به مخلوط‌ها تعمیم می‌دهد. \lr{MLE} استفاده می‌شود اما برای حل تکرارشونده به الگوریتم \lr{EM} نیاز دارد.
    \item \textbf{مدل‌های مارکوف پنهان (\lr{HMM}):} پارامترها با استفاده از \lr{MLE} از طریق الگوریتم جلو-عقب تخمین زده می‌شوند.
\end{itemize}

\section*{۸. خلاصه}
\begin{itemize}
    \item \textbf{درست‌نمایی در مقابل احتمال:} ما درست‌نمایی داده‌ها به شرط پارامترها را بیشینه می‌کنیم، نه احتمال پارامترها به شرط داده‌ها را.
    \item \textbf{لگاریتم درست‌نمایی:} انتقال به حوزه لگاریتم حاصل‌ضرب‌ها را به جمع تبدیل کرده و بهینه‌سازی را ساده می‌کند.
    \item \textbf{نتایج تک‌متغیره:} $\hat{\mu} = \text{mean}(D)$ و $\hat{\sigma}^2 = \text{var}(D)$ (دارای انحراف).
\end{itemize}

\end{document}
